\section*{Open Questions}
\addcontentsline{toc}{section}{Open Questions}

Five substantive open questions grounded in a re-read of the paper (Li~\emph{et~al.} 2012)
and its cited neighbourhood. Each is a genuine unresolved item---not a checklist
item---and carries a concrete falsifiable next-step programme.

\subsection*{Q1. Higher-order secondary jumps ($L^\ast/m$ for $m\geq 3$)}

\textbf{Basis.} REPASS-1 C9 resolved the $L^\ast/2$ jump cleanly (2.24$\times$ drop at
$\bar L=23$\,bp, within 1\,bp of the prediction) but stopped at $\bar L=15$\,bp = $L^\ast/3$
= the $L_m$ dropout threshold. Paper mentions $L^\ast/m$ for arbitrary integer $m$ in a
single sentence on page~4 without deriving jump amplitudes. Whether the $m{=}3$ plateau
exists (masked by $L_m$ cutoff) or has vanishingly small amplitude is unresolved.

\textbf{Next steps.} (1) Lower $L_m$ to 5\,bp (well below biology) and re-scan
$\bar L \in [5,15]$\,bp with 300 runs/point; look for a 3rd plateau. (2) Derive
expected jump amplitude analytically from the recruit-join-release event graph;
hypothesis: jump $n$ has amplitude $\propto \Pr(\text{fragment needs $n$ release events})
\sim (\bar L / L^\ast)^n$. (3) Cross-check by exchanging exponential fragment-length
distribution for uniform and gamma; if jumps are distribution-independent, they are
intrinsic; if they shift, the exponential draw is masking finer structure.

\subsection*{Q2. LET-dependence of the recruitment rate $k_1$}

\textbf{Basis.} Paper treats $k_1, k_2, k_3$ as LET-independent constants and derives the
LET signature purely from fragment-length statistics. Wet-lab evidence (L\"obrich,
Rothkamm) suggests Ku recruitment kinetics may themselves depend on damage-cluster
complexity, which correlates with LET. If real $k_1$ drops for complex clustered
lesions, the biphasic effect could be enhanced (compounding) or suppressed (masking).

\textbf{Next steps.} (1) Literature survey: extract Ku70/80 or DNA-PKcs foci-recruitment
half-lives vs.\ LET from Reindl~2015, Costes~2007, Asaithamby~2011. (2) Extend the
Gillespie model with $k_1(\text{LET}) = k_{1,0}/(1 + \alpha \cdot \text{LET}/\text{LET}_{\text{ref}})$.
(3) Re-run the C5 high/low-LET comparison with LET-scaled $k_1$ and quantify
$\tau_s/\tau_f$; publishable if the ratio changes by $>2\times$ relative to the
fixed-$k_1$ result.

\subsection*{Q3. NHEJ vs.\ HR vs.\ MMEJ competition}

\textbf{Basis.} Paper's $L_m=15$\,bp gating implicitly encodes ``too short for NHEJ.''
But HR and MMEJ have very different length dependencies (HR: $\sim$50--100\,bp; MMEJ:
5--25\,bp). A single-channel Gillespie model with only one $L_m$ collapses this into a
step function; a three-channel model with different length gates would predict a very
different mis-rejoining spectrum.

\textbf{Next steps.} (1) Add HR (rate $k_{HR}$, length gate 50\,bp) and MMEJ (rate
$k_{MMEJ}$, length gate 5--25\,bp with mis-rejoin bias). (2) Calibrate rate ratios from
cell-cycle-phase data (HR active in S/G2). (3) Predict aberration spectra and compare
to Loucas~2013 or Cornforth-Bedford. Falsifiable: MMEJ dominance in the 15--45\,bp
window would flatten the $L^\ast$ jump.

\subsection*{Q4. Diffusion-limited spatial geometry}

\textbf{Basis.} Paper assumes well-mixed fragments in nuclear volume $V$; joining rate
$k_2/V$ per unordered pair with no spatial locality. In reality DSB endpoints diffuse
$\sim 0.5\,\mu\text{m}^2/\text{s}$ and chromosome territories confine motion to
sub-micron domains (Meaburn-Misteli). Under diffusion-limited pairing, the $L^\ast$
threshold---which arises from Ku loading capacity, not diffusion---could remain sharp
OR be washed out by longer effective search times.

\textbf{Next steps.} (1) Reimplement with a spatially-structured joining kernel
$p_{\text{join}}(d) = \exp(-d/d_0)$, $d_0 \sim 0.3\text{--}0.5\,\mu$m. (2) Sample
fragment positions from uniform-in-nucleus or Rabl chromosome territories. (3) Measure
whether the sharp $L^\ast$ jump persists, softens, or migrates to a new
$L^\ast_{\text{eff}}$. Falsifiable: if the jump is absent in the diffusion-limited
version, the paper's mean-time landscape is a well-mixed artifact and the biological
prediction weakens substantially.

\subsection*{Q5. Chromosome aberration dose-response (dicentrics, translocations)}

\textbf{Basis.} Paper produces rejoining-time distributions but does NOT score
chromosome aberrations. The biological punchline (mis-rejoining $\to$ aberrations $\to$
cell death) is asserted structurally (C8) but never quantified. Real datasets
(Cornforth-Bedford dicentric curves, Loucas~2013 mFISH data) show characteristic
linear-quadratic dose-response and specific aberration-type ratios that a
two-fragment Gillespie model should predict without additional free parameters.

\textbf{Next steps.} (1) Extend fragment representation to (length, chromosome, arm,
position) with initial condition = Poisson DSBs per chromosome scaled by track
structure. (2) Score aberrations: cross-chromosome joins $\to$ translocations;
flipped-orientation same-chromosome joins $\to$ inversions; unrejoined ends $\to$
lethal. (3) Fit $D \mapsto f_{\text{dicentric}}$ and compare to Cornforth-Bedford.
Publishable either way: kinetics naturally reproducing the linear-quadratic validates
the mechanism; failure indicates missing physics (track-structure spatial correlations).
